\chapter{Literature Review}
\label{ch:literature}

\section{BasePaper I: Habib and Pradel (2022)}
\label{sec:bp1}

\textbf{S.\ Habib and M.\ Pradel, ``Static Code Analysis Alarms Filtering Reloaded:
A New Real-World Dataset and its ML-Based Utilization,''
\textit{Empirical Software Engineering}, 2022.} \cite{habib2022sca}

This paper is the primary methodological reference for the present project.
The authors identify three pervasive limitations in prior SCA false-positive
filtering research: (i)~approaches target only a small subset of warning types;
(ii) evaluations use synthetic or small benchmarks; (iii) datasets are closed and
non-replicable. They address all three by mining GitHub history to construct a
dataset of \textbf{224,484 real-world SonarQube warnings} from 9,958 Java
projects, labelled using explicit developer actions (fix commit = true positive;
\texttt{//NOSONAR} = false positive). The data covers 160 SonarQube rule checks.

The classification pipeline consists of:
\begin{enumerate}
  \item \textbf{Feature extraction}: the warning line and $n$~surrounding lines
        are tokenized using the \texttt{javalang} lexer.
  \item \textbf{Code embedding}: tokens are embedded with word2vec
        (64-dimensional vectors, window~5) trained on 0.5M Java files.
  \item \textbf{Classification}: four ML classifiers are evaluated ---
        Decision Tree (DTree), Naive Bayes (NBayes), Random Forest (RForest),
        Neural Network (NeuralNet) --- using 10-fold cross-validation with
        stratified sampling per SonarQube warning type and upsampling for
        class imbalance.
\end{enumerate}

\begin{figure}[htbp]
  \centering
  \includegraphics[width=0.9\textwidth]{Figures/fig_bp1_pipeline.png}
  \caption{Classification pipeline of BasePaper~I \cite{habib2022sca}. Real-world
           SonarQube warnings are labelled via developer actions, embedded with
           word2vec, and classified by an ensemble model.}
  \label{fig:bp1_pipeline}
\end{figure}

\textbf{Key results:} Random Forest achieves accuracy~=~0.91, precision~=~0.874,
recall~=~0.760, F1~=~0.813, AUC~=~0.953 on 224,484 real-world warnings across
160 SonarQube rule types. The model correctly filters out 76\% of false positive
warnings while misclassifying only 8\% of true positives as false, preserving
developer attention for real issues.

\textbf{Relevance to this project:} This work directly motivates the use of a
Random Forest classifier for SARIF alert triage in our pipeline. The code-context
embedding strategy and word2vec representation are adapted for our feature set.

\section{BasePaper II: Syed et al.\ (2025)}
\label{sec:bp2}

\textbf{T.A.\ Syed et al., ``Agentic AI for Autonomous Defense in Software Supply
Chain Security: Beyond Provenance to Vulnerability Mitigation,'' arXiv:2512.23480,
2025.} \cite{syed2025agentic}

Syed et al.\ present a multi-agent agentic AI framework for autonomous software
supply chain security. Five specialised agents (Code Analysis, Dependency
Intelligence, CI/CD Monitoring, Access Control, Configuration Audit) are
orchestrated by LangChain and LangGraph. LLM semantic reasoning and PPO-based RL
policies drive mitigation decisions; all agent actions are recorded in a
permissioned blockchain ledger; CI/CD integration is achieved via the Model Context
Protocol (MCP). The framework achieves F1-scores above 0.93 across four
vulnerability classes and reduces mean time-to-mitigation from 28~min to
$<$6~min compared to rule-based systems.

\textbf{Relevance:} Provides the agentic AI design principles (perceive--reason
--act--learn loop, LangGraph conditional routing, LLM chain-of-thought reasoning)
that inform the agentic triage and remediation layer of the present system.

\section{Third Reference: Charoenwet et al.\ (2026)}
\label{sec:3rd}

\textbf{W.\ Charoenwet et al., ``AgenticSCR: Autonomous Secure Code Review for
Immature Vulnerabilities Detection,'' \textit{FSE}, 2026.}
\cite{charoenwet2026agenticscr}

AgenticSCR is an agentic framework for pre-commit vulnerability detection using a
detector--validator subagent pair. The detector subagent loads SAST rules as
long-term semantic memory; the validator cross-checks findings against a CWE
knowledge tree. Evaluated on SCRBench (144 pre-commit changes, 33 CWE types),
it achieves 17.5\% overall correct comments --- 153\% higher than a static LLM
baseline --- while generating 2--5$\times$ fewer comments.

\textbf{Relevance:} The detector--validator pattern motivates the pipeline design
where the ML classifier (detector) and agentic triage layer (validator) are
distinct but complementary stages. The use of SAST rules as semantic memory
informs feature design.

\section{Additional References}
\label{sec:additional}

The following works provide further grounding for the techniques employed.

\begin{description}
  \item[\textbf{Breiman (2001) \cite{breiman2001randomforests}}]
    Introduced the Random Forest algorithm, the foundational ensemble method
    used in our ML triage layer. The parallel bagging approach underpins the
    model's robustness to noisy, high-dimensional feature vectors.

  \item[\textbf{OWASP Benchmark \cite{owasp2023benchmark}}]
    The primary evaluation dataset containing 2,766 labelled Java test cases
    across eight vulnerability categories, enabling ground-truth evaluation
    of SAST and ML classifiers.

  \item[\textbf{GitHub Security Lab / CodeQL \cite{githubseclab2023codeql}}]
    CodeQL serves as the SAST engine generating the SARIF alert reports that
    feed directly into the ML feature extraction pipeline.

  \item[\textbf{Allamanis et al.\ (2018) \cite{allamanis2018mlbigcode}}]
    A comprehensive survey of machine learning techniques applied to large
    code corpora, informing the feature engineering and code embedding
    strategy used in this project.

  \item[\textbf{Iranmanesh et al.\ (2025) \cite{iranmanesh2025zerofalse}}]
    Proposes ZeroFalse, an LLM-guided precision improvement layer on top of
    static analysis --- closely related to our ML triage step but using
    prompting rather than trained classifiers.

  \item[\textbf{Shrestha et al.\ (2025) \cite{shrestha2025agentic}}]
    Explores agentic AI for autonomous software supply chain security,
    providing complementary perspectives on autonomous security pipelines.

  \item[\textbf{FinTech SSRN (2025) \cite{fintech2025vuln}}]
    Demonstrates AI-enhanced combined static and dynamic code analysis in a
    real-world financial technology context, validating the practical impact
    of hybrid security pipelines.

  \item[\textbf{Sooksomstarn et al.\ (2025) \cite{sooksomstarn2025agentic}}]
    Presents agentic hyper-intelligence for cyber threat detection at the
    network and application level, extending the agentic paradigm beyond
    code review into runtime defense.
\end{description}

\section{Synthesis}

Table~\ref{tab:litreview} positions the primary references in terms of their
stage in the vulnerability management lifecycle. Together they form a coherent
rationale for the present project: ML-based alert filtering (BasePaper~I)
augmented by agentic AI reasoning following the principles of BasePaper~II
and the additional supporting works.

\begin{table}[htbp]
  \centering
  \caption{Summary and Positioning of Primary References}
  \label{tab:litreview}
  \renewcommand{\arraystretch}{1.35}
  \begin{tabular}{p{3cm}p{2.5cm}p{4cm}p{2.5cm}}
    \toprule
    \textbf{Reference} & \textbf{Stage} &
    \textbf{Key Contribution} & \textbf{Gap Remaining} \\
    \midrule
    Habib \& Pradel \cite{habib2022sca} &
      Post-scan alert triage &
      224K real-world labelled SonarQube warnings; RForest 91\% accuracy &
      Static dataset; no agentic remediation \\
    Syed et al.\ \cite{syed2025agentic} &
      Supply chain defense &
      Multi-agent LLM+RL framework with MCP and blockchain &
      No ML-based alert filtering \\
    Charoenwet et al.\ \cite{charoenwet2026agenticscr} &
      Pre-commit review &
      Detector--validator subagents; 153\% improvement over static LLM &
      No ML triage \\
    Iranmanesh et al.\ \cite{iranmanesh2025zerofalse} &
      FP reduction &
      LLM-based precision improvement on static alerts &
      No autonomous remediation \\
    \bottomrule
  \end{tabular}
\end{table}
